\documentclass[11pt,letterpaper]{article}
\usepackage[T1]{fontenc}
\usepackage{amsmath,amsthm,mathtools}
\usepackage{newpxtext,newpxmath}
\usepackage[margin=1in,headheight=15pt]{geometry}
\usepackage{microtype}
\usepackage{xcolor,booktabs,array,tabularx,enumitem}
\usepackage[most]{tcolorbox}
\usepackage{fancyhdr,titlesec,needspace,etoolbox}
\usepackage{xurl}
\usepackage{hyperref,bookmark}
\definecolor{Accent}{HTML}{24596C}
\definecolor{Ink}{HTML}{20313A}
\definecolor{Pale}{HTML}{F1F6F7}
\hypersetup{colorlinks=true,linkcolor=Accent,citecolor=Accent,urlcolor=Accent,
 pdftitle={Cofinality and Scalar Extension in Surcomplex Entire Function Theory},
 pdfauthor={OpenAI ChatGPT},
 pdfsubject={Hahn series, arbitrary-rank valuations, entire functions, cofinality, units, divisors of zeros}}
\titleformat{\section}{\Large\bfseries\color{Accent}}{\thesection}{0.7em}{}
\titleformat{\subsection}{\large\bfseries\color{Ink}}{\thesubsection}{0.7em}{}
\pagestyle{fancy}\fancyhf{}
\fancyhead[L]{\small\scshape Cofinality and surcomplex entire functions}
\fancyhead[R]{\small Research article}
\fancyfoot[C]{\thepage}
\renewcommand{\headrulewidth}{0.35pt}
\setlength{\parskip}{0.3em}
\setlength{\parindent}{1.2em}
\setlist{topsep=4pt,itemsep=3pt,parsep=0pt,leftmargin=1.6em}
\setcounter{tocdepth}{2}
\pretocmd{\section}{\Needspace{7\baselineskip}}{}{}
\pretocmd{\subsection}{\Needspace{5\baselineskip}}{}{}
\numberwithin{equation}{section}
\newtheorem{theorem}{Theorem}[section]
\newtheorem{lemma}[theorem]{Lemma}
\newtheorem{proposition}[theorem]{Proposition}
\newtheorem{corollary}[theorem]{Corollary}
\theoremstyle{definition}
\newtheorem{definition}[theorem]{Definition}
\newtheorem{example}[theorem]{Example}
\theoremstyle{remark}
\newtheorem{remark}[theorem]{Remark}
\newcommand{\No}{\mathbf{No}}
\newcommand{\C}{\mathbb C}
\newcommand{\R}{\mathbb R}
\newcommand{\Q}{\mathbb Q}
\newcommand{\N}{\mathbb N}
\newcommand{\Z}{\mathbb Z}
\newcommand{\K}{K_\Gamma}
\newcommand{\KD}{K_\Delta}
\newcommand{\HH}{\mathcal H_\Gamma}
\newcommand{\TT}{\mathcal T_\Gamma}
\newcommand{\EE}{\mathcal E_\Gamma}
\newcommand{\OO}{\mathcal O_\Gamma}
\newcommand{\mm}{\mathfrak m_\Gamma}
\newcommand{\supp}{\operatorname{supp}}
\newcommand{\cf}{\operatorname{cf}}
\newcommand{\ord}{\operatorname{ord}}
\newcommand{\lc}{\operatorname{lc}}
\newcommand{\ini}{\operatorname{in}}
\newcommand{\Div}{\operatorname{Div}}
\newcommand{\dom}{\operatorname{Dom}}
\newcommand{\red}[1]{\overline{#1}}
\newcommand{\repocommit}{aa846271b4dcae2c055b216126a87210292ec19b}
\newcommand{\code}[1]{\texttt{#1}}
\newtcolorbox{scopebox}{colback=Pale,colframe=Accent,boxrule=0.5pt,
 arc=1mm,left=9pt,right=9pt,top=7pt,bottom=7pt,breakable}
\newcolumntype{Y}{>{\raggedright\arraybackslash}X}

\begin{document}
\hypersetup{pageanchor=false}
\begin{titlepage}
\vspace*{5mm}
{\color{Accent}\large\scshape Surreal and surcomplex research\par}
\vspace{10mm}
{\color{Ink}\Huge\bfseries Cofinality and\par Scalar Extension in\par Surcomplex Entire\par Function Theory\par}
\vspace{6mm}
{\Large Sharp convergence and unit criteria\par beyond rank one\par}
\vspace{9mm}
{\large Prepared for Vladimir Reshetnikov\par}
{\large OpenAI ChatGPT\quad\textbullet\quad 21 September 2026\par}
\vspace{8mm}
\begin{abstract}
Let $K_\Gamma=\C((t^\Gamma))$ be a Hahn field with a nonzero divisible ordered
abelian value group, canonically viewed as a surcomplex workspace when
$\Gamma\subset\No$. We study ordinary power series whose evaluations are
strongly Hahn-summable at every point of this field. We prove that this
apparently support-theoretic condition is equivalent to
$v(a_n)/n\to+\infty$ in the whole ordered group. Consequently nonpolynomial
entire series exist exactly when $\Gamma$ has countable cofinality.

A different invariant governs local inversion: a nonconstant valuation-restricted
series is a unit precisely when its constant coefficient strictly dominates
and the least positive valuation gap is an order unit of $\Gamma$.
For every ordered-group extension $\Gamma\subset\Delta$, we then determine
the exact common strong evaluation domain of all nonpolynomial entire series
over $K_\Gamma$. It is the valuation ring associated with the convex hull of
$\Gamma$ in $\Delta$. Entireness survives exactly for cofinal extensions,
and no new zeros appear anywhere in that domain. Support-controlled
preparation supplies the necessary arbitrary-rank root theory and a genus-zero
factorization theorem. Explicit infinite-rank surreal examples separate
countable cofinality, topological nilpotence, strong summation, and extension
to all surcomplex numbers.
\end{abstract}
\vfill
\begin{scopebox}\small
\textbf{Research status.} The main arbitrary-rank classifications and their
combined scalar-extension formulation are proposed research contributions;
no matching statements were located in the primary sources reviewed here.
This is not a priority certificate or a claimed solution of a named published
conjecture. Classical Hahn algebra and rank-one function theory are explicitly
credited. The proofs are supplied in full apart from named standard inputs;
they have not been independently refereed or formally verified. The accompanying
program checks 664 finite exact assertions, not the infinite theorems.
\end{scopebox}
\end{titlepage}
\hypersetup{pageanchor=true}
\pagenumbering{roman}
\tableofcontents
\clearpage
\pagenumbering{arabic}

\section{The questions and the contributions}\label{sec:introduction}

\subsection{Why the value group is part of the analytic object}
The phrase ``an entire surcomplex function'' hides a choice of both a summation
rule and an allowed collection of scales. In a fixed Hahn field, a series can
be evaluable everywhere without being a polynomial. Enlarging the field can
make that very same defining series meaningless at new points. Moreover,
in higher rank a positive valuation need not give a sequence tending to zero.
These are structural issues, not defects cured by choosing a sufficiently
small ordinary real number.

The Surreal repository already contains a rank-one treatment over
$\C((t^\R))$, including preparation, zero counting, zero-free rigidity and
canonical products. Its foundations discussion distinguishes fixed-workspace
entireness from assertions on all of $\No[i]$ \cite{RepoMap,RepoRankOne}.
The task here is not to repeat those results with new terminology. We ask for
\emph{sharp arbitrary-rank criteria}: exactly which ordered groups support
nonpolynomial entire series, exactly which restricted series are invertible,
and exactly which points remain available after scalar extension.

Throughout the article, a power series has ordinary nonnegative integer
degrees. Its coefficient supports may have arbitrary set-sized ordinal order
types. The two kinds of indexing must not be conflated.

\subsection{The principal results}
Fix a nonzero divisible ordered abelian group $\Gamma$ and set
\[
 K_\Gamma=\C((t^\Gamma)),\qquad
 v\left(\sum_\gamma c_\gamma t^\gamma\right)=\min\supp(c),\qquad v(0)=+\infty.
\]
Precise definitions of strong summability and the algebras occur in
Sections~\ref{sec:support}--\ref{sec:cofinality}. The results can already be
stated informatively.

\begin{scopebox}
\textbf{I. The cofinality theorem.}
For $f=\sum a_nZ^n$, strong evaluation at every point of $K_\Gamma$ is equivalent to
\[
 \frac{v(a_n)}{n}\longrightarrow+\infty\quad(n\ge1)
\]
in the full order of $\Gamma$, with zero coefficients assigned value $+\infty$.
Thus the entire algebra contains a nonpolynomial if and only if
$\cf(\Gamma)=\aleph_0$.
\end{scopebox}

This is stronger than a sufficient growth estimate. It shows that global
strong summability automatically recovers intrinsic valuation convergence,
even though the two notions differ on a single disk.

\begin{scopebox}
\textbf{II. The restricted-unit theorem.}
Let $f=\sum a_nZ^n$ be nonconstant with $v(a_n)\to+\infty$. Then $f$ has an
inverse of the same restricted kind if and only if $a_0\ne0$ and
\[
 \delta=\min_{n\ge1,\ a_n\ne0}v(a_n/a_0)>0,
 \qquad \{m\delta:m\in\N\}\text{ is cofinal in }\Gamma.
\]
\end{scopebox}

The second condition is essential. A positive $\delta$ with cofinal multiples
is called an \emph{order unit}. Countable cofinality does not imply the
existence of an order unit. Consequently there are fields supporting abundant
nonpolynomial entire functions whose restricted algebra has no nonconstant
units at all.

\begin{scopebox}
\textbf{III. The universal scalar-extension theorem.}
Let $\Gamma\subset\Delta$ be divisible ordered groups and let $f$ be any
nonpolynomial entire series over $K_\Gamma$. Its strong evaluation domain
in $K_\Delta$ is exactly
\[
 \mathcal D_{\Gamma,\Delta}
 =\{0\}\cup\{z\ne0:v_\Delta(z)\ge\gamma
                       \text{ for some }\gamma\in\Gamma\}.
\]
This domain is independent of $f$. It is a coarsened valuation ring.
The series remains entire precisely when $\Gamma$ is cofinal in $\Delta$.
Every zero in this domain already belongs to $K_\Gamma$, with unchanged
multiplicity. When the extension is not cofinal, there is not even a different
entire series over $K_\Delta$ agreeing with $f$ on $K_\Gamma$.
\end{scopebox}

The exact domain statement is not merely a radius-of-convergence analogy.
Outside the indicated ring the leading exponents contain a strictly descending
subsequence, so strong summation is impossible. Inside the ring, a support
argument works even when the original sequence of partial sums no longer
converges in the enlarged valuation topology.

The fourth major result is a divisor classification and genus-zero product
formula valid in arbitrary rank (Section~\ref{sec:divisors}). It supports the
main conclusions, but its rank-one antecedents are classical and are not
claimed as new.

\subsection{What was already known, and what is being proposed}
Hahn-field operations, positive-support geometric inversion, and algebraic
closedness for divisible value group and algebraically closed coefficient field
are established results \cite{Poonen,Neumann}. Surreal normal forms identify
these fields with concrete subfields of $\No[i]$ when the exponent group is a
subgroup of $\No$ \cite{Gonshor,BournezGuilmant}.

In rank one, the convergence criterion, preparation, Newton polygons, and
canonical products belong to standard non-Archimedean function theory.
Cherry's lectures give the one-variable theory, and his factorization paper
explicitly presents its factorization results as known \cite{CherryLectures,CherryFactorization}.
The distinction between positive valuation and topological nilpotence in
higher rank is also known: Conrad gives an explicit rank-two example in his
notes on Huber rings \cite[Section~6.2]{Conrad}.

The proposed contribution is the exact combination of \emph{global strong}
summability, arbitrary ordered-group cofinality, the necessary-and-sufficient
restricted-unit test, and the universal domain and zero-conservation theorem
under arbitrary ordered-group extensions. The reviewed sources do not state
this package or the infinite-rank examples below. This is evidence of a useful
research direction and of a distinction from the repository's existing work,
not proof that no equivalent theorem exists elsewhere.

The repository audit was targeted: its catalogue, relevant introductory
material, and the rank-one entire-function and product sections were inspected
at snapshot
\begin{center}\small\texttt{\repocommit}.\end{center}
It was not a line-by-line certification of every source manuscript. No theorem
below depends for its validity on the correctness of an unrefereed repository
report.

\section{Hahn workspaces and summation with support control}\label{sec:support}

\subsection{Set-sized foundations and surreal interpretation}
An ordered group in this article is a \emph{set}, and all supports and indexing
families used in a Hahn sum are sets. Ordinary ZFC therefore suffices for the
field-theoretic arguments. The optional interpretation in $\No$ can be made
with definable classes, or in a class theory, without treating the whole
surreal class as a set.

For $\Gamma\subset\No$, our convention is
\begin{equation}\label{eq:monomial}
 t^\gamma:=\omega^{-\gamma}.
\end{equation}
The right side is Conway's monomial map, not an instruction to apply a general
surreal exponential with base $\omega$. The normal-form theorem identifies
$\R((t^\Gamma))$ with an ordered subfield of $\No$, and coefficientwise
complexification gives $K_\Gamma\subset\No[i]$. Complex conjugation acts only
on coefficients. None of the proofs uses the Berarducci--Mantova derivation;
all derivatives below differentiate the coordinate $Z$ and fix $K_\Gamma$.

Write
\[
 \mathcal O_\Gamma=\{a:v(a)\ge0\},\qquad
 \mathfrak m_\Gamma=\{a:v(a)>0\}.
\]
The residue map $\mathcal O_\Gamma\to\C$ is extraction of the exponent-zero
coefficient. The residue field is $\C$, so every nonzero integer has valuation
zero. Divisibility of $\Gamma$ ensures that $K_\Gamma$ is algebraically closed
\cite[Corollary~4]{Poonen}; this is the principal algebraic input to zero counting.

\subsection{Strong summability}
\begin{definition}\label{def:strong}
A family $(x_i)_{i\in I}$ of Hahn series is \emph{strongly summable} if
\[
 S=\bigcup_{i\in I}\supp(x_i)
\]
is well ordered and, for every $\gamma\in\Gamma$, only finitely many $i$ have
a nonzero coefficient at $t^\gamma$. Its sum is obtained by the resulting
finite coefficientwise additions. Zero members cause no difficulty.
\end{definition}

This is a condition on the \emph{family before cancellation}. A displayed
infinite expression cannot be declared summable merely because a speculative
cancellation would leave a small answer.

We use two standard support facts. If $A,B$ are well-ordered subsets of an
ordered group, then $A+B$ is well ordered and each $\gamma$ has only finitely
many decompositions $a+b=\gamma$ with $a\in A,b\in B$. If $S\subset\Gamma_{>0}$
is well ordered, the additive monoid
\[
 S^*=\{s_1+\cdots+s_m:m\ge0,\ s_i\in S\}
\]
is well ordered and a fixed group element has only finitely many
representations as a finite word in $S$, over all lengths. These are the
Hahn--Neumann support lemmas \cite{Neumann,Poonen}. In particular, for a
positive-support Hahn element $h$, the family $(h^m)_{m\ge0}$ is strongly
summable and $\sum_{m\ge0}h^m=(1-h)^{-1}$.

The finite-word assertion is important: finiteness separately at each word
length would not justify summing over all lengths. The lemmas hold for
arbitrary ordered abelian groups, not just subgroups of $\R$.

\begin{lemma}[Cofinal valuations imply strong summability]\label{lem:cofinal-strong}
Let $(x_n)_{n\ge0}$ be Hahn series such that $v(x_n)\to+\infty$ in $\Gamma$.
Then the family is strongly summable and its partial sums converge to its
Hahn sum in the intrinsic valuation topology.
\end{lemma}
\begin{proof}
For each $\beta\in\Gamma$, only finitely many $x_n$ have support meeting
$(-\infty,\beta]$. Thus a fixed exponent appears only finitely often. To see
that the union is well ordered, take a nonempty subset of it and choose one
of its elements $\beta$. All its elements at most $\beta$ lie in the finite
union of the supports just described. That finite union is well ordered and
therefore supplies a least element of the original subset. Finally every
tail sum has valuation at least the least valuation of its summands. Eventually
all those valuations exceed any prescribed $\beta$, which proves convergence.
\end{proof}

The converse fails. In a rank-two group with $0<\epsilon\ll\eta$, the family
$(t^{n\epsilon})_{n\ge0}$ is strongly summable, but its valuations never exceed
$\eta$. Its partial sums do not converge in the full intrinsic valuation
topology. This distinction will be used, not suppressed.

\subsection{The polynomial-coefficient Hahn algebra}
Define
\begin{equation}\label{eq:H}
 \mathcal H_\Gamma=\C[Z]((t^\Gamma)).
\end{equation}
Thus $F\in\mathcal H_\Gamma$ means
\[
 F=\sum_{\gamma\in S}P_\gamma(Z)t^\gamma,
 \qquad S\subset\Gamma\text{ well ordered},\quad P_\gamma\in\C[Z].
\]
There is no common bound on the degrees of the $P_\gamma$. Regrouping by $Z$
gives an injective ring map $\mathcal H_\Gamma\hookrightarrow K_\Gamma[[Z]]$.
A series $f=\sum a_nZ^n$ belongs to $\mathcal H_\Gamma$ exactly when the
coefficient family $(a_n)_n$ is strongly summable. In particular,
\begin{equation}\label{eq:rescale-H}
 (a_nt^{ns})_n\text{ is strongly summable}
 \quad\Longleftrightarrow\quad f(t^sZ)\in\mathcal H_\Gamma.
\end{equation}
The finite degree of each $P_\gamma$ is precisely the pointwise finiteness
condition at exponent $\gamma$.

Let $w(F)$ be its least exponent, and let $\mathcal H_{\Gamma,\ge0}$ and
$\mathcal H_{\Gamma,>0}$ be the nonnegative and positive support parts.
The latter is an ideal of the former. Their quotient is $\C[Z]$.
By the support lemma,
\begin{equation}\label{eq:strong-unit}
 U\in1+\mathcal H_{\Gamma,>0}
 \quad\Longrightarrow\quad U^{-1}\in1+\mathcal H_{\Gamma,>0}.
\end{equation}
This is a strong Hahn inverse. It is not automatically a valuation-restricted
power series.

\begin{lemma}[Evaluation on integral points, including after extension]
\label{lem:integral-eval}
Suppose $\Gamma\subset\Delta$ are ordered groups. Every
$F\in\mathcal H_{\Gamma,\ge0}$ can be strongly evaluated at every
$u\in\mathcal O_\Delta$. Evaluation is a ring homomorphism into
$\mathcal O_\Delta$, and positive-support elements evaluate into
$\mathfrak m_\Delta$. Consequently $1+\mathcal H_{\Gamma,>0}$ evaluates to
nonzero elements at every such point. The same evaluation is defined for
all $F\in\mathcal H_\Gamma$ after a scalar monomial normalization.
\end{lemma}
\begin{proof}
Write $u=c+\varepsilon$, with $c\in\C$ and either $\varepsilon=0$ or
$v(\varepsilon)>0$. Each polynomial $P_\gamma(c+\varepsilon)$ is a finite
Taylor expansion. If $T=\supp(\varepsilon)$, all output exponents lie in
$S+T^*$. The support lemmas make this set well ordered and give finitely many
pairs $(\gamma,\tau)$ with $\gamma+\tau$ fixed. For each such pair there are
finitely many words producing $\tau$, and the polynomial at $\gamma$ has
finite degree. All coefficient manipulations are therefore finite.
This argument also justifies distributing products and regrouping, proving
the ring-homomorphism assertion.

All exponents are nonnegative. If the original support is strictly positive,
adding exponents from $T^*$ leaves it strictly positive. An evaluated element
of $1+\mathcal H_{\Gamma,>0}$ thus has residue $1$. Finally every nonzero
$F\in\mathcal H_\Gamma$ has a least exponent $\mu$; applying the preceding
argument to $t^{-\mu}F$ proves the last assertion.
\end{proof}

This lemma proves strong summability of the original degree-indexed evaluation
family as well: before regrouping, the same finite polynomial Taylor
expansions and the same finite support decompositions account for every
coefficient. It is stronger than merely assigning a value after a formal
rearrangement.

\section{Global strong entireness and the cofinality dichotomy}
\label{sec:cofinality}

\subsection{The exact convergence criterion}
\begin{definition}
The \emph{strong entire algebra} $\mathcal E_\Gamma$ consists of formal power
series $f=\sum_{n\ge0}a_nZ^n\in K_\Gamma[[Z]]$ such that
$(a_nz^n)_{n\ge0}$ is strongly summable for every $z\in K_\Gamma$.
The resulting evaluation is denoted $f(z)$. Its algebra structure will follow
from the convergence criterion.
\end{definition}

In statements about limits, $b_n\to+\infty$ means that for every
$\beta\in\Gamma$ there exists $N$ such that $b_n>\beta$ for all $n\ge N$.
This quantifies over \emph{all} ranks of the group.

\begin{theorem}[Global strong convergence criterion]\label{thm:criterion}
For $f=\sum_{n\ge0}a_nZ^n\in K_\Gamma[[Z]]$, the following are equivalent.
\begin{enumerate}[label=(\roman*)]
\item $f\in\mathcal E_\Gamma$.
\item $(a_nt^{ns})_n$ is strongly summable for every $s\in\Gamma$.
\item $v(a_n)+ns\to+\infty$ for every $s\in\Gamma$.
\item $v(a_n)/n\to+\infty$ for $n\ge1$.
\end{enumerate}
Zero coefficients have value $+\infty$ in these formulas. Every evaluation
of an entire series is also the intrinsic valuation limit of its partial sums.
\end{theorem}
\begin{proof}
The implication (i)$\Rightarrow$(ii) uses the points $z=t^s$.
To prove (ii)$\Rightarrow$(iii), fix $s$ and put
$b_n=v(a_n)+ns$ for the nonzero coefficients. Strong summability at $t^s$
implies that the $b_n$ have a common lower bound $\alpha$. If they do not tend
to $+\infty$, there exist $\beta$ and an increasing sequence of indices
$n_1<n_2<\cdots$ with
\[
 \alpha\le b_{n_j}\le\beta.
\]
Choose $\delta>\max(\beta-\alpha,0)$. Then
\begin{align*}
 &(b_{n_{j+1}}-n_{j+1}\delta)-(b_{n_j}-n_j\delta)\\
 &\hspace{12mm}\le\beta-\alpha-(n_{j+1}-n_j)\delta<0.
\end{align*}
These are leading exponents of summands evaluated at $t^{s-\delta}$. They form
a strictly decreasing sequence, contradicting the well-ordered support
required by (ii). Therefore (iii) holds.

Assume (iii). Given $\theta\in\Gamma$, set $s=-\theta$. Eventually
$v(a_n)-n\theta>0$, so $v(a_n)/n>\theta$. This proves (iv).
Conversely, assume (iv), and fix $s,\beta\in\Gamma$. Choose
$\theta>-s+\max(\beta,0)$. Eventually
\[
 v(a_n)+ns>n(\theta+s)>\beta,
\]
which is (iii). Divisibility is used only to write the equivalent slope
condition without passing to a divisible hull.

Finally, for $z\ne0$ with $v(z)=s$, the summand valuation is exactly
$v(a_n)+ns$. Lemma~\ref{lem:cofinal-strong} proves strong summability and
valuation convergence. At $z=0$ only the constant term remains. Thus
(iii)$\Rightarrow$(i), and the last assertion follows as well.
\end{proof}

The decisive step is to move to the slightly larger monomial scale
$t^{s-\delta}$. Strong summability at one scale can hide bounded positive
valuation progress. Requiring every scale exposes it as a decreasing support.

\subsection{Classification of the permitted value groups}
\begin{definition}
The cofinality $\cf(\Gamma)$ is the least cardinality of a cofinal subset of
its ordered underlying set. A nonzero ordered group has no largest element,
so its cofinality is infinite.
\end{definition}

\begin{theorem}[Cofinality dichotomy]\label{thm:cofinality}
The following conditions are equivalent:
\[
 \mathcal E_\Gamma\ne K_\Gamma[Z],
 \qquad \mathcal E_\Gamma\text{ contains a nonpolynomial},
 \qquad \cf(\Gamma)=\aleph_0.
\]
In particular, if $\cf(\Gamma)>\aleph_0$, every strong entire ordinary power
series is a polynomial.
\end{theorem}
\begin{proof}
If $f$ has infinitely many nonzero coefficients, Theorem~\ref{thm:criterion}
says their slopes $v(a_n)/n$ eventually exceed every member of $\Gamma$.
They form a countable cofinal subset.

Conversely, choose an increasing positive cofinal sequence
$\gamma_1<\gamma_2<\cdots$ in $\Gamma$. Then
\begin{equation}\label{eq:constructed-entire}
 f(Z)=1+\sum_{n\ge1}t^{n\gamma_n}Z^n
\end{equation}
has slope sequence $\gamma_n\to+\infty$ and is nonpolynomial. The criterion
proves that it is entire.
\end{proof}

The conclusion is about cofinality, not cardinality or rank. An uncountable
rank-one group such as $\R$ has countable cofinality. A group of infinite rank
can have countable cofinality as well. Conversely, a sufficiently long
well-ordered hierarchy of Archimedean classes produces uncountable
cofinality and polynomial rigidity.

\begin{proposition}[Algebra operations]\label{prop:algebra}
$\mathcal E_\Gamma$ is a $K_\Gamma$-algebra. For every fixed $s$, the partial
sums converge in the weighted Gauss valuation
\[
 w_s(f)=\min_n\bigl(v(a_n)+ns\bigr)
\]
when $f\ne0$. Formal differentiation and termwise integration, with any
chosen constant of integration, preserve $\mathcal E_\Gamma$.
\end{proposition}
\begin{proof}
For sums the weighted valuations of coefficients are bounded below by the
minimum of the corresponding two weighted valuations. For products, set
$A_i=v(a_i)+is$ and $B_j=v(b_j)+js$. Both sequences tend to infinity and
have finite lower bounds $A,B$. Given $\beta$, choose $N$ so that
$A_i>\beta-B$ and $B_i>\beta-A$ for $i>N$. If $n>2N$, each pair $i+j=n$
has $i>N$ or $j>N$, so the corresponding product coefficient has weighted
valuation greater than $\beta$. The criterion applies to the convolution.
The same tail estimate gives weighted convergence of partial sums.

Differentiation replaces $a_n$ by $na_n$ and lowers the degree by one;
integration replaces it by $a_n/(n+1)$ and raises the degree by one.
Nonzero integers have valuation zero. At each fixed $s$ these operations only
shift the weighted bounds by $-s$ or $s$.
\end{proof}

\section{Restricted series: the sharp unit criterion}\label{sec:units}

\subsection{An algebra different from the strong disk algebra}
Define the valuation-restricted algebra
\begin{equation}\label{eq:T}
 \mathcal T_\Gamma
 =\left\{\sum_{n\ge0}a_nZ^n:v(a_n)\to+\infty\text{ in }\Gamma\right\}.
\end{equation}
Then
\[
 K_\Gamma[Z]\subset\mathcal E_\Gamma\subset\mathcal T_\Gamma
 \subset\mathcal H_\Gamma\subset K_\Gamma[[Z]].
\]
The first two inclusions need not be strict. The last two capture different
summation requirements. The product proof of Proposition~\ref{prop:algebra},
with $s=0$, proves that $\mathcal T_\Gamma$ is a ring.

For $0\ne F\in\mathcal H_\Gamma$, let
\[
 w_0(F)=\min_n v(a_n),\qquad
 \ini_0(F)=\red{t^{-w_0(F)}F}\in\C[Z]\setminus\{0\}.
\]
The Hahn leading-coefficient calculation gives
\begin{equation}\label{eq:gauss-multiplicative}
 w_0(FG)=w_0(F)+w_0(G),\qquad
 \ini_0(FG)=\ini_0(F)\ini_0(G).
\end{equation}
In particular this holds in $\mathcal T_\Gamma$.

\begin{definition}
An element $\delta>0$ is an \emph{order unit} of $\Gamma$ if for every
$\gamma\in\Gamma$ some positive integer $m$ satisfies $m\delta>\gamma$.
Equivalently, $m\delta\to+\infty$. Equivalently, the monomial $t^\delta$
is topologically nilpotent in the intrinsic valuation topology of $K_\Gamma$.
\end{definition}

The order-unit condition is stronger than positivity. Its familiar topological
meaning and the failure of positivity to suffice in rank two are discussed in
\cite[Section~6.2]{Conrad}. The following theorem identifies its exact role
in the coefficient algebra \eqref{eq:T}.

\subsection{Necessity and sufficiency}
\begin{theorem}[Restricted-unit criterion]\label{thm:units}
Let $f=\sum a_nZ^n\in\mathcal T_\Gamma$ be nonconstant. It is a unit in
$\mathcal T_\Gamma$ if and only if $a_0\ne0$ and
\begin{equation}\label{eq:unit-gap}
 \delta(f)=\min_{n\ge1,\ a_n\ne0}v(a_n/a_0)
\end{equation}
is a positive order unit of $\Gamma$. Nonzero constant series are always
units. The minimum in \eqref{eq:unit-gap} exists whenever $a_0\ne0$.
\end{theorem}
\begin{proof}
The minimum exists because the coefficient valuations tend to infinity: below
any chosen nonzero coefficient's valuation there are only finitely many
nonzero coefficients.

Suppose $fg=1$ with $g\in\mathcal T_\Gamma$. The constant coefficient forces
$a_0\ne0$. Divide $f$ by $a_0$ and multiply $g$ by $a_0$, so that both
constant coefficients are $1$. Then $w_0(f),w_0(g)\le0$, while
\eqref{eq:gauss-multiplicative} gives their sum as zero. Hence both are zero.
Their reductions in $\C[Z]$ multiply to $1$, so both reductions are constant.
Every nonconstant coefficient of $f$ and of $g$ consequently has positive
valuation. In particular $\delta=\delta(f)>0$.

Assume for contradiction that $\delta$ is not an order unit. Its generated
convex subgroup
\begin{equation}\label{eq:Hdelta}
 H=\{\gamma\in\Gamma:|\gamma|\le m\delta
                       \text{ for some }m\in\N\}
\end{equation}
is proper. Coarsen the valuation to the ordered quotient $\Gamma/H$.
Every coefficient of $f$ and $g$ is integral for this coarsening. Their
coarsened reductions are \emph{polynomials} over the coarsened residue field:
choose $\beta>H$, meaning $\beta>h$ for every $h\in H$. Since coefficient
valuations tend to infinity in $\Gamma$, only finitely many coefficients
can have value at most $\beta$. All remaining coefficients reduce to zero.

The reduced polynomial of $f$ is nonconstant. Indeed a nonconstant coefficient
with valuation $\delta\in H$ has nonzero coarsened residue. On the other hand,
the reduced polynomials of $f$ and $g$ multiply to $1$ in a polynomial ring
over a field. A nonconstant polynomial cannot be a unit. This contradiction
proves that $\delta$ is an order unit.

Conversely, suppose $a_0\ne0$ and $\delta(f)$ is a positive order unit. Write
$f=a_0(1+h)$, where $h$ has zero constant coefficient and $w_0(h)=\delta$.
The formal inverse is
\begin{equation}\label{eq:restricted-inverse}
 f^{-1}=a_0^{-1}\sum_{m\ge0}(-h)^m.
\end{equation}
Its summands have Gauss valuation $m\delta\to+\infty$. The sum exists strongly
in $\mathcal H_\Gamma$. To prove that it belongs to $\mathcal T_\Gamma$,
fix a valuation threshold $\beta$. All sufficiently large $m$ contribute
only coefficients of valuation greater than $\beta$. The remaining finitely
many powers of $h$ belong to $\mathcal T_\Gamma$, so their coefficients
also eventually exceed $\beta$. Thus the coefficients of
\eqref{eq:restricted-inverse} tend to infinity. The finite geometric identity,
followed by the justified strong sum, proves $ff^{-1}=1$.
\end{proof}

The coarsening argument explains why cancellations in a hypothetical inverse
cannot evade the order-unit obstruction. It is a necessity theorem, not merely
an observation that one particular geometric-series construction fails.

\subsection{Consequences and a two-scale counterexample}
\begin{corollary}\label{cor:unit-structure}
The algebra $\mathcal T_\Gamma$ has a nonconstant unit if and only if $\Gamma$
has an order unit. If it has none, then
$\mathcal T_\Gamma^\times=K_\Gamma^\times$.
The rule ``constant coefficient strictly dominates all the others'' is a
sufficient unit criterion for every restricted series if and only if $\Gamma$
is Archimedean.
\end{corollary}
\begin{proof}
An order unit $\delta$ makes $1-t^\delta Z$ a nonconstant unit by the theorem.
The converse follows by applying the theorem to any nonconstant unit.
An ordered group is Archimedean exactly when every positive element has
cofinal positive integer multiples. Apply this characterization to
$1-t^\delta Z$ for the last assertion.
\end{proof}

\begin{example}[Zero-free does not mean restricted-invertible]
\label{ex:ranktwo-unit}
Let $\Gamma=\Q\epsilon\oplus\Q\eta$ with
$0<\epsilon\ll\eta$; the $\eta$ coordinate is compared first. The polynomial
\[
 f(Z)=1-t^\epsilon Z
\]
has residue $1$ and no zero in $\mathcal O_\Gamma$: its unique zero is
$t^{-\epsilon}$, of negative valuation. Yet its only formal inverse,
\[
 \sum_{n\ge0}t^{n\epsilon}Z^n,
\]
is not restricted, since $n\epsilon<\eta$ for every $n$. It does belong to
$\mathcal H_\Gamma$ and can be strongly evaluated on the integral disk.
In contrast, $1-t^\eta Z$ has a restricted inverse because $n\eta$ is cofinal.
\end{example}

Thus one must not import a rank-one Banach proof with the phrase
``the valuation increases by a positive amount at each step'' left unchanged.
For a valid valuation-convergence proof that positive amount must be an
order unit. The next section instead uses positive \emph{supports} and
finite-word control; it works without any order unit.

\section{Preparation without a rank-one contraction}\label{sec:preparation}

\subsection{A support-certified factorization}
The preparation result below is a Hahn-support version of a familiar
Weierstrass mechanism. Its role is to supply a correct proof in arbitrary
rank, rather than to claim a new rank-one preparation theorem. The coefficient
ring is always $\mathcal H_\Gamma$, not silently $\mathcal T_\Gamma$.

\begin{theorem}[Polynomial preparation in the strong disk algebra]
\label{thm:preparation}
Let
\[
 F=P_0+E\in\mathcal H_{\Gamma,\ge0},\qquad
 P_0\in\C[Z]\text{ monic of degree }d,\qquad
 E\in\mathcal H_{\Gamma,>0}.
\]
There is a unique factorization
\begin{equation}\label{eq:preparation}
 F=P U,
 \qquad P=P_0+R,\quad \deg R<d,
 \qquad U=1+V,
\end{equation}
where every coefficient of $R$ has positive valuation and
$V\in\mathcal H_{\Gamma,>0}$. The polynomial $P$ is monic of degree $d$.
If $S=\supp(E)$, all exponents of $R$ and $V$ belong to
$S^*\setminus\{0\}$. The construction is compatible with ordered-group
extensions.
\end{theorem}
\begin{proof}
If $E=0$, take $R=V=0$. Otherwise perform Euclidean division by $P_0$
\emph{coefficientwise} in $\C[Z]$. This produces two $\C$-linear operators:
quotient $q$ and remainder $r$, with
\[
 H=P_0q(H)+r(H),\qquad \deg_Zr(H)<d.
\]
These operators do not introduce new Hahn exponents.

We construct terms homogeneous in the number of occurrences of the
perturbation $E$. Put $H_1=E$ and
\[
 V_1=q(H_1),\qquad R_1=r(H_1).
\]
For $j\ge2$, having constructed the smaller indices, put
\begin{equation}\label{eq:prep-recursion}
 H_j=-\sum_{\substack{a+b=j\\a,b\ge1}}R_aV_b,
 \qquad V_j=q(H_j),\qquad R_j=r(H_j).
\end{equation}
Induction shows that the support of each $R_j,V_j$ is contained in the
$j$-fold sum $jS$. The positive-support lemma makes the families
$(R_j)_j,(V_j)_j$ jointly summable across all $j$: a fixed exponent is produced
by only finitely many finite words in $S$, and the recursion at each finite
word length uses finitely many algebraic operations.

Define $R=\sum_{j\ge1}R_j$ and $V=\sum_{j\ge1}V_j$. Their coefficients have
positive valuation, and $\deg R<d$ because every remainder does.
The equations for $j=1$ and $j\ge2$, with the justified regrouping of double
sums, give
\[
 E=R+P_0V+RV.
\]
This is exactly $(P_0+R)(1+V)=P_0+E$.

For uniqueness, suppose $(P_0+R)(1+V)=(P_0+R')(1+V')$. If a difference is
nonzero, let $\lambda>0$ be the least exponent occurring in $R-R'$ or
$V-V'$. The coefficient at $t^\lambda$ in the difference of the two products
is
\[
 (R-R')_\lambda+P_0(V-V')_\lambda=0.
\]
All cross terms with a positive-support factor occur above $\lambda$.
The first summand has degree less than $d$; the second is divisible by the
monic polynomial $P_0$. Uniqueness of polynomial division makes both zero,
contradicting the choice of $\lambda$.

The operations in \eqref{eq:prep-recursion} preserve exponents and use only
polynomial arithmetic over $\C$. The same support certificate remains valid
under an ordered embedding into a larger group, proving compatibility.
\end{proof}

The use of perturbation \emph{length} in the proof is bookkeeping for a Hahn
sum. It does not assert that the error tends to zero as the length tends to
infinity in the full valuation topology. That assertion would generally be
false when the smallest exponent of $E$ is not an order unit.

\subsection{Zeros on the integral disk}
\begin{corollary}[Finite zero geometry]\label{cor:prepared-zeros}
With the hypotheses of Theorem~\ref{thm:preparation}, the zeros of $F$ in
$\mathcal O_\Gamma$ are exactly the roots of $P$, counted with multiplicity.
There are exactly $d$ of them. All roots of $P$ lie in $\mathcal O_\Gamma$.
For each $c\in\C$, the number with residue $c$ equals the multiplicity of
$c$ as a root of $P_0$.
\end{corollary}
\begin{proof}
The coefficients of $P$ are integral, and $P$ is monic. If a root $r$ had
$v(r)<0$, its leading term $r^d$ would have strictly smaller valuation than
every lower-degree term, and their sum could not be zero. Thus all its roots
are integral. There are $d$ roots with multiplicity because $K_\Gamma$ is
algebraically closed. Lemma~\ref{lem:integral-eval} gives $U(u)\ne0$ for all
integral $u$, so $F$ has exactly those roots.

Reducing the split factorization
$P(Z)=\prod_{j=1}^d(Z-r_j)$ gives
$P_0(Z)=\prod_{j=1}^d(Z-\red{r_j})$, proving the residue count.
Finally polynomial substitution $Z=u+W$ at an integral $u$ is legitimate in
$\mathcal H_\Gamma$ by the same support argument as
Lemma~\ref{lem:integral-eval}. Its unit factor has nonzero constant term,
so the order of vanishing of $F$ equals that of $P$.
\end{proof}

Every nonzero $F\in\mathcal H_\Gamma$ can be brought into the theorem's form:
if $\mu=w_0(F)$ and $c$ is the leading coefficient of its initial polynomial,
then
\[
 (ct^\mu)^{-1}F=P_0+E
\]
with $P_0$ monic and $E$ of positive support. In particular a nonzero element
of $\mathcal H_\Gamma$ has only finitely many zeros on the integral disk.
This finite-zero statement does not require ordinary complex compactness.

\subsection{A concrete preparation}
Let $\epsilon>0$ be arbitrary, even if it is not an order unit, and abbreviate
$T=t^\epsilon$. Consider
\begin{equation}\label{eq:cubic}
 F(Z)=Z^2-Z+T(Z^3+1).
\end{equation}
The first terms of the recursive preparation are
\begin{align}
 P(Z)&=Z^2-Z+T(Z+1)-T^2(3Z+1)+T^3(8Z+2)+\cdots,
 \label{eq:prep-example-P}\\
 U(Z)&=1+T(Z+1)-T^2+3T^3+\cdots.
 \label{eq:prep-example-U}
\end{align}
Here the ellipses mean terms supported on $m\epsilon$ with $m\ge4$;
they do not assert valuation convergence. The polynomial $P$ has two
integral roots, with residues $0$ and $1$. The cubic $F$ has one further root
outside the integral disk. Its valuation is $-\epsilon$: the sum of the three
roots is $-T^{-1}$, and the other two roots are integral, so the third has
that valuation.

The companion program verifies \eqref{eq:preparation} through perturbation
length eight for this example. The infinite identity follows from the support
proof, not from extrapolating the finite test.

\section{Newton geometry and elementary entire-function algebra}
\label{sec:zeros}

\subsection{Initial polynomials at arbitrary scales}
For a nonzero $f=\sum a_nZ^n\in\mathcal E_\Gamma$ and $s\in\Gamma$, define
\begin{align}
 \mu_s(f)&=\min_n\bigl(v(a_n)+ns\bigr),\label{eq:mu}\\
 \ini_s(f)(X)&=\red{t^{-\mu_s(f)}f(t^sX)}\in\C[X]\setminus\{0\},
 \label{eq:initial}\\
 d_s(f)&=\deg\ini_s(f),\qquad
 e_s(f)=\ord_{X=0}\ini_s(f).\label{eq:de}
\end{align}
The minimum is attained, and only finitely many indices attain it, because
the weighted coefficient valuations tend to infinity. The active indices
are exactly the nonzero coefficients of the initial polynomial.

\begin{theorem}[Root counts and residue directions]\label{thm:root-counts}
For $0\ne f\in\mathcal E_\Gamma$, let zeros always be counted with
multiplicity. At every $s\in\Gamma$,
\begin{align*}
 \#\{z:f(z)=0,\ v(z)\ge s\}&=d_s(f),\\
 \#\{z:f(z)=0,\ v(z)>s\}&=e_s(f),\\
 \#\{z:f(z)=0,\ v(z)=s\}&=d_s(f)-e_s(f).
\end{align*}
For $c\in\C$, the number of roots in
$\{z:v(z)\ge s,\ \red{t^{-s}z}=c\}$ equals the multiplicity of $c$ in
$\ini_s(f)$.
\end{theorem}
\begin{proof}
Normalize $f(t^sX)$ by its least monomial and the leading complex coefficient
of its initial polynomial. Apply Theorem~\ref{thm:preparation} and
Corollary~\ref{cor:prepared-zeros}. The coordinate change $z=t^sX$
identifies the integral disk with the closed ball $v(z)\ge s$. The zero
residue direction is the open ball, so it contributes the order $e_s(f)$.
The other residue directions form the shell.
\end{proof}

\begin{proposition}[Locally finite Newton profile]\label{prop:newton}
On each ordered interval $[s_0,s_1]\subset\Gamma$, the function
$s\mapsto\mu_s(f)$ is the minimum of a \emph{finite} set of affine functions
$v(a_n)+ns$. All changes of active indices occur at slopes
\[
 s=\frac{v(a_i)-v(a_j)}{j-i}\in\Gamma\qquad(i\ne j).
\]
If $g\in\mathcal E_\Gamma$ satisfies $\mu_s(g)>\mu_s(f)$, then $f+g$ has the
same initial polynomial and all the same counts and residue directions at
scale $s$.
\end{proposition}
\begin{proof}
Choose any index $j$ with $a_j\ne0$ and set $M=v(a_j)+js_1$. If $n$ is
active at some $s\in[s_0,s_1]$, then
\[
 v(a_n)+ns_0\le v(a_n)+ns=\mu_s(f)\le v(a_j)+js\le M.
\]
Only finitely many $n$ satisfy this inequality by
Theorem~\ref{thm:criterion} at $s_0$. Two different affine functions can
be equal only at the displayed slope. Finally the strict perturbation
inequality leaves every coefficient at the least Hahn exponent unchanged.
\end{proof}

This supplies the expected Newton and Rouch\'e conclusions in the chosen
category. Their rank-one predecessors are classical \cite{CherryLectures};
the point here is that preparation and the finite active-index argument did
not require a real-valued norm.

\subsection{Translation and division by a zero}
\begin{proposition}[Entire Taylor translation]\label{prop:translation}
For $f\in\mathcal E_\Gamma$ and $a\in K_\Gamma$, the series
\begin{equation}\label{eq:translation}
 f(a+Z)=\sum_{k\ge0}b_kZ^k,\qquad
 b_k=\sum_{n\ge k}\binom nk a_n a^{n-k},
\end{equation}
is well defined and belongs to $\mathcal E_\Gamma$. It evaluates to the
indicated composition at every point, and
$b_k=f^{(k)}(a)/k!$.
\end{proposition}
\begin{proof}
The case $a=0$ is immediate. Otherwise put $\rho=v(a)$ and fix $s\in\Gamma$.
With $t_0=\min(s,\rho)$, the valuation of the summand contributing to
$b_kt^{ks}$ is at least
\[
 v(a_n)+(n-k)\rho+ks\ge v(a_n)+nt_0.
\]
The right side tends cofinally to infinity with $n$. For each $n$ there are
only finitely many $k$, so the entire triangular family of binomial summands
is strongly summable. Regrouping is legitimate. For $k\to\infty$, every
contributing index has $n\ge k$, so the same tail estimate proves
$v(b_k)+ks\to+\infty$. The criterion gives entireness. The binomial identity
and the corresponding justified substitution prove the evaluation formula;
formal differentiation gives the coefficient identity.
\end{proof}

The order of vanishing at $a$ can consequently be defined as the least $k$
with $b_k\ne0$. It agrees with the local polynomial multiplicity in
Theorem~\ref{thm:root-counts}. This also makes its independence from the
chosen containing ball explicit.

\begin{lemma}[Division by a linear zero factor]\label{lem:linear-division}
If $f\in\mathcal E_\Gamma$ and $f(r)=0$, there is $g\in\mathcal E_\Gamma$
with $f(Z)=(Z-r)g(Z)$.
\end{lemma}
\begin{proof}
For $r=0$, shift the coefficients. For $r\ne0$, define
\begin{equation}\label{eq:linear-division}
 b_n=\sum_{k\ge n+1}a_kr^{k-n-1},\qquad g(Z)=\sum_{n\ge0}b_nZ^n.
\end{equation}
Each inner family has valuations tending to infinity. Fix $s$, put
$\rho=v(r)$ and $t_0=\min(s,\rho)$. Its summands after weighting by $ns$
have valuation at least
\[
 v(a_k)+(k-n-1)\rho+ns\ge v(a_k)+kt_0-t_0.
\]
Since $k\ge n+1$, these lower bounds tend to infinity uniformly as
$n\to\infty$. Hence $g$ is entire. Subtracting the adjacent sums gives
$b_{n-1}-rb_n=a_n$ for $n\ge1$. The constant identity
$-rb_0=a_0$ is exactly $f(r)=0$. Thus $(Z-r)g=f$.
\end{proof}

\subsection{Rigidity and value distribution}
\begin{theorem}[Zero-free rigidity and finite-zero classification]
\label{thm:zero-free}
A zero-free member of $\mathcal E_\Gamma$ is a nonzero constant.
A nonzero member has finitely many zeros if and only if it is a polynomial.
\end{theorem}
\begin{proof}
If $a_0=0$ then zero is a root. Otherwise suppose $a_n\ne0$ for some $n>0$.
Choose $s$ so small that $v(a_n)+ns<v(a_0)$. The initial polynomial at that
scale has positive degree, and Theorem~\ref{thm:root-counts} gives a root.
Therefore a zero-free series is constant.

If $f$ has finitely many zeros, divide out all their linear factors with their
multiplicities using Lemma~\ref{lem:linear-division}. The resulting entire
series has no zeros and is constant, so $f$ is a polynomial. The converse is
ordinary polynomial algebra over the algebraically closed field $K_\Gamma$.
\end{proof}

\begin{corollary}\label{cor:fibers}
Every nonconstant entire function $K_\Gamma\to K_\Gamma$ is surjective.
Every nonpolynomial entire function has infinitely many points in each fiber.
Every injective entire function is affine of degree one.
\end{corollary}
\begin{proof}
For any $c\in K_\Gamma$, the function $f-c$ is nonconstant and cannot be
zero-free. If $f$ is nonpolynomial, so is $f-c$; it therefore has infinitely
many zeros. An injective entire function must consequently be a polynomial.
In characteristic zero, a polynomial of degree $d>1$ has $d$ distinct roots
in a fiber over any value outside its finite set of critical values. It
cannot be injective. Affine polynomials of degree one are injective.
\end{proof}

These value-distribution conclusions have familiar non-Archimedean forms;
they are included to describe the resulting theory, not offered as independent
priority claims. The arbitrary-rank novelty lies in the hypotheses under which
this theory exists and survives enlargement.

\section{The universal domain under scalar extension}\label{sec:extension}

\subsection{The convex-hull valuation ring}
Let $\Gamma\subset\Delta$ be an ordered inclusion of nonzero divisible
ordered abelian groups, with the canonical inclusion
$K_\Gamma\subset K_\Delta$. Define the convex hull
\begin{equation}\label{eq:convex-hull}
 H=\operatorname{conv}_\Delta(\Gamma)
 =\{\delta\in\Delta:|\delta|\le\gamma
                      \text{ for some }\gamma\in\Gamma_{>0}\}.
\end{equation}
It is a divisible convex subgroup of $\Delta$. Coarsening $v_\Delta$ by $H$
gives a valuation with value group $\Delta/H$.

\begin{lemma}\label{lem:coarsened-ring}
The valuation ring of this coarsening is
\begin{equation}\label{eq:D}
 \mathcal D_{\Gamma,\Delta}
 =\{0\}\cup\{z\ne0:v_\Delta(z)\ge\gamma
                       \text{ for some }\gamma\in\Gamma\}.
\end{equation}
Its maximal ideal consists of zero and the elements satisfying
$v_\Delta(z)>\gamma$ for every $\gamma\in\Gamma$.
Its residue field is canonically $\C((t^H))$. Moreover,
$\mathcal D_{\Gamma,\Delta}=K_\Delta$ exactly when $\Gamma$ is cofinal in
$\Delta$.
\end{lemma}
\begin{proof}
A value $s$ has nonnegative class modulo $H$ exactly when it is bounded below
by some member of $H$. Every member of $H$ is itself bounded below by a
member of $\Gamma$, proving \eqref{eq:D}. Its class is strictly positive
exactly when $s>H$, equivalently $s>\Gamma$.

For an element of this valuation ring, discard all support exponents outside
$H$. No negative coset can occur, and the discarded positive cosets form the
coarsened maximal ideal. The remaining Hahn series has support in $H$.
This is a surjective residue homomorphism to $\C((t^H))$ with the stated
kernel. Finally the ring is the whole field exactly when $H=\Delta$, which
by \eqref{eq:convex-hull} is equivalent to cofinality of $\Gamma$ in $\Delta$.
\end{proof}

The residue field here is generally not $\C$. It records all scales in the
convex hull $H$, whereas the original fine valuation's residue field records
only the exponent zero.

\subsection{Exact strong evaluation, not just a sufficient domain}
\begin{theorem}[Universal scalar-extension domain]\label{thm:extension}
Let $f=\sum a_nZ^n\in\mathcal E_\Gamma$ be nonpolynomial. In $K_\Delta$,
the family $(a_nz^n)_n$ is strongly summable if and only if
$z\in\mathcal D_{\Gamma,\Delta}$. Thus its domain is independent of the
choice of nonpolynomial $f$. Its values on this domain again lie in
$\mathcal D_{\Gamma,\Delta}$.
\end{theorem}
\begin{proof}
Suppose $z\ne0$ lies in the stated ring. Choose $\gamma\in\Gamma$ with
$v_\Delta(z)\ge\gamma$ and put $u=t^{-\gamma}z\in\mathcal O_\Delta$.
Since $f$ is entire over $K_\Gamma$,
\[
 f(t^\gamma X)\in\mathcal H_\Gamma.
\]
Lemma~\ref{lem:integral-eval} gives strong summability of its evaluation
family at $u$, which is the original family $(a_nz^n)_n$. If
$\mu=\mu_\gamma(f)\in\Gamma$, the same lemma applied after normalization
gives $v_\Delta(f(z))\ge\mu$ unless the value is zero. Hence the value lies
in $\mathcal D_{\Gamma,\Delta}$. At $z=0$ the assertion is immediate.

Conversely, suppose $s=v_\Delta(z)<\gamma$ for every $\gamma\in\Gamma$.
List the infinitely many nonzero coefficients at increasing indices
$n_1<n_2<\cdots$ and write $\alpha_j=v(a_{n_j})\in\Gamma$. Since $s<0$,
\begin{align*}
 (\alpha_{j+1}+n_{j+1}s)-(\alpha_j+n_js)
 &=\alpha_{j+1}-\alpha_j+(n_{j+1}-n_j)s\\
 &\le\alpha_{j+1}-\alpha_j+s<0.
\end{align*}
The final inequality uses $s<\alpha_j-\alpha_{j+1}$, a member of $\Gamma$.
These are actual leading exponents of the nonzero summands $a_{n_j}z^{n_j}$,
regardless of the higher terms of $z$. Their strict descent prevents the
support union from being well ordered. Thus the family is not strongly
summable.
\end{proof}

\begin{corollary}[Cofinal extensions are exactly the entire extensions]
\label{cor:entire-extension}
For nonpolynomial $f\in\mathcal E_\Gamma$, its same coefficient series
belongs to $\mathcal E_\Delta$ if and only if $\Gamma$ is cofinal in $\Delta$.
Polynomial series remain entire under every such extension.
\end{corollary}
\begin{proof}
Combine Theorem~\ref{thm:extension} with
Lemma~\ref{lem:coarsened-ring}. Polynomial evaluations are finite.
\end{proof}

A cofinal extension need not preserve rank, and a noncofinal extension need
not introduce any algebraic element: the obstruction is entirely in the
ordered hierarchy of scales. Nor does a strongly defined value in the larger
field have to be a limit there. For example, if $\Delta$ contains a positive
value above all of $\Gamma$, the original cofinal coefficient valuations in
$\Gamma$ need not tend to infinity in $\Delta$.

\subsection{No new zeros, even for an infinite divisor}
\begin{theorem}[Zero conservation under arbitrary enlargement]
\label{thm:zero-conservation}
Let $0\ne f\in\mathcal E_\Gamma$ and let $\Gamma\subset\Delta$ be as above.
Every zero of its strongly defined extension on
$\mathcal D_{\Gamma,\Delta}$ belongs to $K_\Gamma$. Its multiplicity is
unchanged. This remains true when $f$ has infinitely many zeros.
\end{theorem}
\begin{proof}
Suppose $z$ lies in the domain and $f(z)=0$. Choose $\gamma\in\Gamma$
with $v_\Delta(z)\ge\gamma$. Preparation of $f(t^\gamma X)$ over
$K_\Gamma$ gives
\[
 f(t^\gamma X)=ct^\mu P_\gamma(X)U_\gamma(X),
 \qquad U_\gamma\in1+\mathcal H_{\Gamma,>0},
\]
with $P_\gamma$ monic over $K_\Gamma$. Its unit factor stays nonzero at
every $X\in\mathcal O_\Delta$. Therefore
$P_\gamma(t^{-\gamma}z)=0$. The polynomial already splits into linear
factors over the algebraically closed field $K_\Gamma$, so
$t^{-\gamma}z$, and hence $z$, lies in $K_\Gamma$.

The same prepared polynomial has the same linear factor multiplicities in
both fields. The evaluated unit is nonzero in both, so local orders of
vanishing are unchanged. The argument uses only the finitely many zeros
in one containing ball and applies separately to every zero; it does not
require that the global divisor be finite.
\end{proof}

\subsection{No alternative entire continuation}
\begin{proposition}[Identity on constants]\label{prop:identity}
Suppose two formal series over $K_\Delta$ belong to
$\mathcal H_\Delta$ and agree at infinitely many distinct points of
$\mathcal O_\Delta$. Then they are equal as formal series.
In particular, agreement at all ordinary complex constants suffices.
\end{proposition}
\begin{proof}
Their difference belongs to $\mathcal H_\Delta$. If it were nonzero,
preparation would give only finitely many zeros in the integral disk.
\end{proof}

\begin{theorem}[Obstruction to any entire-series continuation]
\label{thm:no-continuation}
Suppose $\Gamma$ is not cofinal in $\Delta$ and
$f\in\mathcal E_\Gamma$ is nonpolynomial. There is no
$g\in\mathcal E_\Delta$ agreeing with $f$ at all points of $K_\Gamma$.
In fact, there is no such $g$ agreeing with $f$ at all points of $\C$.
\end{theorem}
\begin{proof}
The coefficient series of $f$ belongs to $\mathcal H_\Delta$, since its
strong support certificate at scale zero over $\Gamma$ remains valid under
an ordered embedding. The same is true of $g$. Agreement on $\C$ would force
the two coefficient series to be identical by Proposition~\ref{prop:identity}.
But Corollary~\ref{cor:entire-extension} says that this coefficient series
is not entire over $K_\Delta$, a contradiction.
\end{proof}

This rules out an alternative entire \emph{power-series} continuation.
It does not rule out arbitrary functions extending the pointwise values,
or continuations in a different analytic category. No density of $\C$ in
a fine surreal topology is being asserted or used.

\section{Divisors and genus-zero factorization in arbitrary rank}
\label{sec:divisors}

\subsection{The correct finiteness condition}
An effective divisor here is a set-sized multiset of points of $K_\Gamma$,
with nonnegative integer multiplicities. Call it \emph{ball-finite} if
\begin{equation}\label{eq:ball-finite}
 \sum_{v(z)\ge s}m(z)<\infty\qquad\text{for every }s\in\Gamma.
\end{equation}
This is stronger and more useful than merely saying that the points are
isolated in a topology. It is exactly the property supplied by
Theorem~\ref{thm:root-counts} for the zeros of a nonzero entire series.

\begin{lemma}[Automatic countability and escape]\label{lem:divisor-countable}
Every ball-finite effective divisor is finite or countably infinite.
In the infinite case its nonzero points, repeated with multiplicity, can be
enumerated as $(r_j)_{j\ge1}$ with
\[
 -v(r_j)\longrightarrow+\infty\quad\text{in }\Gamma.
\]
In particular an infinite ball-finite divisor forces
$\cf(\Gamma)=\aleph_0$.
\end{lemma}
\begin{proof}
First choose countably many occurrences from an infinite divisor. Finiteness
in every ball says that only finitely many of the chosen occurrences have
$-v(r_j)\le\beta$, for any $\beta\in\Gamma$. Thus their negative valuations
tend to infinity and form a countable cofinal set. The corresponding
countable collection of closed balls covers $K_\Gamma$. Each ball contains
only finitely many occurrences of the whole divisor, so the whole divisor is
countable. Any enumeration without repeating occurrences then has the stated
escape property. The point zero has finite multiplicity by
\eqref{eq:ball-finite} and is separated off.
\end{proof}

The countability conclusion has not been assumed from a real-valued absolute
value or a countable family of radii. The divisor itself produces the
countable cofinal family of radii. When $\Gamma$ has uncountable cofinality,
all ball-finite divisors are necessarily finite.

\subsection{Constructing the product with no exponential factors}
\begin{theorem}[Prescribed zeros]\label{thm:prescribed-divisor}
Let $D$ be a ball-finite effective divisor, with multiplicity $m$ at zero and
nonzero occurrences $(r_j)$. The product
\begin{equation}\label{eq:canonical-product}
 G_D(Z)=Z^m\prod_j\left(1-\frac{Z}{r_j}\right)
\end{equation}
defines an element of $\mathcal E_\Gamma$. Its zeros are exactly $D$.
The definition is independent of the escaping enumeration. If the product is
infinite, its partial products converge in the weighted Gauss valuation at
every scale $s\in\Gamma$.
\end{theorem}
\begin{proof}
There is nothing to prove for a finite product. In the infinite case use
Lemma~\ref{lem:divisor-countable}. Fix $s\in\Gamma$ and substitute $Z=t^sX$.
Apart from finitely many indices, the terms
\[
 h_j(X)=-\frac{t^sX}{r_j}
\]
have positive Gauss valuation, and their Gauss valuations tend cofinally to
infinity. Their joint positive support is well ordered and locally finite
by Lemma~\ref{lem:cofinal-strong}, applied to their polynomial-coefficient
Hahn series.

Expand the tail product over finite subsets $J$ of the tail index set:
\begin{equation}\label{eq:subset-product}
 \prod_j(1+h_j)=\sum_{J\text{ finite}}\prod_{j\in J}h_j.
\end{equation}
For a fixed output exponent, Neumann's lemma permits only finitely many
finite words in the joint positive support. Each exponent in that support
occurs in only finitely many $h_j$. Hence only finitely many subset products
contribute. The right side is a well-defined element of
$1+\mathcal H_{\Gamma,>0}$. Multiplying by the omitted finite initial product
and by $t^{ms}X^m$ gives an element of $\mathcal H_\Gamma$.

At $s=0$ this constructs the formal coefficients of $G_D$. At every other
$s$ the same finite-subset expressions construct those coefficients multiplied
by $t^{ns}$. Thus $G_D(t^sX)\in\mathcal H_\Gamma$ for every $s$, and
Theorem~\ref{thm:criterion} proves entireness. Reordering does not change the
finite-subset sum.

To count zeros in the ball at scale $s$, separate all factors with
$v(r_j)\ge s$. There are finitely many. Every remaining factor belongs to
a positive-support tail as above, whose product has residue $1$ and is a
unit on the integral disk. Therefore the only zeros in this ball are the
listed finite factors, with their multiplicities. Varying $s$ gives exactly
$D$.

Finally, let $P_N$ be partial products, after including the finitely many
nonpositive-valuation factors at scale $s$. Their remaining tails differ
from $1$ by elements whose least exponent is at least
$\min_{j>N}(s-v(r_j))$, which tends to infinity. Multiplication by the fixed
initial product shifts this lower bound by a fixed value. Hence
$w_s(P_N-G_D)\to+\infty$ (with the harmless $Z^m$ factor included).
\end{proof}

No complex Weierstrass exponential correction factors occur. Here
positive-support products are controlled by finite-word combinatorics, and
the escaping condition also gives genuine intrinsic convergence. The analogous
rank-one product construction is standard \cite{CherryLectures,CherryFactorization};
the proof above explains why it remains valid without a rank-one norm.

\subsection{Uniqueness and the divisor classification}
\begin{theorem}[Divisor classification and canonical factorization]
\label{thm:factorization}
The zero-divisor map induces a bijection
\[
 \bigl(\mathcal E_\Gamma\setminus\{0\}\bigr)/K_\Gamma^\times
 \ \longleftrightarrow\
 \{\text{ball-finite effective divisors on }K_\Gamma\}.
\]
Explicitly, if $f$ has a zero of order $m$ at zero and nonzero zeros $(r_j)$,
repeated with multiplicity, then
\begin{equation}\label{eq:factorization}
 f(Z)=cZ^m\prod_j\left(1-\frac Z{r_j}\right),
 \qquad c=[Z^m]f.
\end{equation}
\end{theorem}
\begin{proof}
Root counting gives ball-finiteness, and
Theorem~\ref{thm:prescribed-divisor} proves existence for every such divisor.
For uniqueness remove the common power $Z^m$ from $f$ and $G_D$; both then
have nonzero constant coefficients, so their quotient $q$ exists formally in
$K_\Gamma[[Z]]$.

At every scale $s$, prepare the two rescaled series. Their monic prepared
polynomials have exactly the same finite multiset of roots in the integral
disk, so they are identical. Cancelling this polynomial leaves a nonzero
scalar times the quotient of two elements of
$1+\mathcal H_{\Gamma,>0}$. That quotient belongs to
$\mathcal H_\Gamma$ by \eqref{eq:strong-unit}. Since the prepared polynomial
has nonzero constant coefficient after removing the common zero at the
origin, cancellation is valid also in the formal power-series field, and
this local quotient is $q(t^sX)$.

Thus $q(t^sX)\in\mathcal H_\Gamma$ for every $s$. Theorem~\ref{thm:criterion}
makes $q$ entire. On every containing ball it is a scalar times a positive-support
unit, so it is zero-free. Theorem~\ref{thm:zero-free} makes it constant.
The constant term after removal of $Z^m$ is $c$, proving
\eqref{eq:factorization} and uniqueness modulo constants.
\end{proof}

\begin{corollary}[Conjugation and real-coefficient functions]
A nonzero entire series has all coefficients in $\R((t^\Gamma))$ if and
only if its zero divisor is invariant under coefficientwise conjugation and
its first nonzero coefficient belongs to $\R((t^\Gamma))$.
\end{corollary}
\begin{proof}
Conjugating coefficients conjugates each value and each zero with its
multiplicity. Conversely, under the stated conditions the conjugate series
has the same divisor and first nonzero coefficient. Uniqueness in
Theorem~\ref{thm:factorization} makes the two series equal.
\end{proof}

\subsection{Recovering the scale boundary from the function}
\begin{corollary}[An intrinsic description of the extension boundary]
\label{cor:intrinsic-boundary}
Let $f\in\mathcal E_\Gamma$ be nonpolynomial. In any larger ordered group
$\Delta$, the following subsets generate the same convex subgroup:
\[
 \Gamma,\qquad
 \{v(a_n)/n:n\ge1,\ a_n\ne0\},\qquad
 \{-v(r):r\ne0,\ f(r)=0\}.
\]
Consequently the universal domain of Theorem~\ref{thm:extension} can be
recovered either from the coefficient slopes or from the zero divisor.
\end{corollary}
\begin{proof}
All the displayed values belong to $\Gamma$. The slopes are cofinal in
$\Gamma$ by Theorem~\ref{thm:criterion}. There are infinitely many zeros by
Theorem~\ref{thm:zero-free}, and their negative valuations are cofinal by
Lemma~\ref{lem:divisor-countable}. A cofinal subset of an ordered group
generates its full convex hull in any containing ordered group.
\end{proof}

This makes the extension obstruction intrinsic. It is not an artifact of
having embedded the coefficients in an unnecessarily large workspace whose
extra scales the function never detects.

\section{Explicit higher-rank and surreal examples}\label{sec:examples}

\subsection{Two series with the same leading valuations and different boundary behavior}
\label{subsec:hidden-tails}
Let $\Gamma=\Q\epsilon\oplus\Q\eta$ with $0<\epsilon\ll\eta$, and define
\begin{align}
 A(Z)&=\sum_{n\ge1}t^{n^2\epsilon}Z^n,\label{eq:hidden-A}\\
 B(Z)&=\sum_{n\ge1}
       \left(t^{n^2\epsilon}+t^{\eta-n^2\epsilon}\right)Z^n.
       \label{eq:hidden-B}
\end{align}
Their nonzero coefficients have exactly the same valuations $n^2\epsilon$.
Nevertheless their \emph{monomial} strong evaluation domains differ.
For $s=a\eta+b\epsilon$ with $a,b\in\Q$,
\begin{equation}\label{eq:hidden-profile}
 \begin{array}{c|ccc}
 &a<0&a=0&a>0\\\hline
 A(t^s)&\text{not strong}&\text{strong}&\text{strong}\\
 B(t^s)&\text{not strong}&\text{not strong}&\text{strong}
 \end{array}
\end{equation}

For $a>0$, the leading valuations have $\eta$ coordinate $na\to+\infty$;
Lemma~\ref{lem:cofinal-strong} proves both positive assertions. For $a<0$,
those leading valuations strictly decrease. For $a=0$, the exponents of
$A(t^s)$ are $(n^2+nb)\epsilon$: after finitely many terms they strictly
increase, with only finite repetition at any exponent. They form a well-ordered
support. The extra exponents of $B(t^s)$ are
\[
 \eta+(nb-n^2)\epsilon,
\]
which strictly decrease for all sufficiently large $n$. They destroy strong
summability despite not affecting a single leading coefficient valuation.

Neither series is entire over the whole rank-two field. Their coefficient
slopes $n\epsilon$ remain below $\eta$, so Theorem~\ref{thm:criterion}
correctly excludes both. There is no contradiction between this example and
the global valuation-only criterion: leading valuations can fail to decide
strong summability at a \emph{boundary scale}, while global summability at
\emph{every} scale removes that ambiguity. Formula~\eqref{eq:hidden-profile}
is stated for monomial arguments; it does not silently assert a classification
for arbitrary nonmonomial arguments of $B$.

\subsection{An infinite-rank entire function with explicit surreal zeros}
\label{subsec:infinite-rank}
Let
\begin{equation}\label{eq:Gamma-infinity}
 \Gamma_\infty=\bigoplus_{j\ge0}\Q e_j,
 \qquad 0<e_0\ll e_1\ll e_2\ll\cdots,
\end{equation}
where the largest index with nonzero coefficient decides the order.
This group has countable cofinality, but no order unit. Indeed every group
element has finite support, and the next basis element exceeds every positive
integer multiple of it.

It is a concrete ordered subgroup of $\No$ by taking
\[
 e_j=\omega^j\qquad(j=0,1,2,\ldots)
\]
and finite rational surreal linear combinations. Consider
\begin{equation}\label{eq:infinite-product}
 P(Z)=\prod_{j\ge0}(1-t^{e_j}Z).
\end{equation}

\begin{theorem}[A genuine infinite-rank surcomplex entire function]
\label{thm:infinite-example}
The product \eqref{eq:infinite-product} belongs to
$\mathcal E_{\Gamma_\infty}$ and is nonpolynomial. Its zeros are exactly
\begin{equation}\label{eq:surreal-zeros}
 r_j=t^{-e_j}=\omega^{\omega^j}\qquad(j\ge0),
\end{equation}
all simple and positive surreal. Its coefficient of $Z^k$ is
\begin{equation}\label{eq:elementary-coeff}
 c_k=(-1)^k\sum_{j_1<\cdots<j_k}t^{e_{j_1}+\cdots+e_{j_k}},
 \qquad
 v(c_k)=e_0+e_1+\cdots+e_{k-1}
\end{equation}
for $k\ge1$, and $c_0=1$.
At scale $s\in\Gamma_\infty$, put
\[
 N_s=\#\{j:e_j+s<0\},\qquad M_s=\#\{j:e_j+s=0\}.
\]
These are finite, $M_s\in\{0,1\}$, and
\begin{align}
 \mu_s(P)&=\sum_{e_j+s<0}(e_j+s),\label{eq:product-profile}\\
 \ini_s(P)(X)&=(-1)^{N_s}X^{N_s}(1-X)^{M_s}.\label{eq:product-initial}
\end{align}
Finally,
\begin{equation}\label{eq:product-derivative}
 v(P'(r_j))=e_j+\sum_{k<j}(e_k-e_j).
\end{equation}
\end{theorem}
\begin{proof}
Since $e_j\to+\infty$, the proposed zero divisor is ball-finite.
Theorem~\ref{thm:prescribed-divisor} proves the product's existence,
entireness, and exact simple zero set. Expanding over finite subsets gives
\eqref{eq:elementary-coeff}. Among the $k$-element subsets, the sum of the
first $k$ basis elements is uniquely least in the largest-index order.
Its coefficient is $(-1)^k$, so it cannot cancel. Every $c_k$ is nonzero.

The slope sequence $(e_0+\cdots+e_{k-1})/k$ is cofinal: once $k-1$ exceeds
the largest index in a prescribed group element, its positive $e_{k-1}$
coordinate dominates that element. This gives a second direct check of
entireness via Theorem~\ref{thm:criterion}.

At scale $s$, only finitely many factors have $e_j+s\le0$. A factor with
negative value contributes that value to $\mu_s$ and contributes $-X$ to the
normalized initial polynomial. A factor of value zero contributes $1-X$;
all positive factors contribute $1$. This proves
\eqref{eq:product-profile}--\eqref{eq:product-initial}.

Differentiate the factor with index $j$ and evaluate at $r_j$. All other
product-rule terms vanish there, giving
\[
 P'(r_j)=-t^{e_j}\prod_{k\ne j}(1-t^{e_k-e_j}).
\]
The factors with $k>j$ form a positive-support unit. The finitely many
factors with $k<j$ have valuations $e_k-e_j$. Adding these valuations proves
\eqref{eq:product-derivative}.
\end{proof}

This example is not confined to any one Archimedean scale. Nevertheless its
entire theory is nontrivial and its zero divisor is completely explicit.
At the same time,
\[
 \mathcal T_{\Gamma_\infty}^{\times}=K_{\Gamma_\infty}^{\times}
\]
by Corollary~\ref{cor:unit-structure}. Thus nonpolynomial entire functions and
nonconstant restricted units are genuinely different existence questions.

\subsection{One extension preserves the function; another destroys entireness}
Adjoin $\omega^{1/2}$ to $\Gamma_\infty$ and take its divisible additive span
with $\Gamma_\infty$. This is an ordered-group extension in which
$\Gamma_\infty$ remains cofinal: a sufficiently large integer power of
$\omega$ dominates every finite combination of the added and old generators.
The same series $P$ remains entire by Corollary~\ref{cor:entire-extension}.

Instead adjoin the new generator
\[
 e_*=\omega^\omega>\Gamma_\infty
\]
and set $\Delta=\Gamma_\infty\oplus\Q e_*$. At the new point
$Z=t^{-e_*}$, the nonzero summands of $P$ have leading exponents
\[
 (e_0+\cdots+e_{k-1})-ke_*.
\]
The difference of consecutive exponents is $e_k-e_*<0$. The defining sum is
therefore not strongly summable. Theorem~\ref{thm:no-continuation} rules out
any alternative entire power series on $K_\Delta$ agreeing with $P$ on the
old field. On the still-defined domain
$\mathcal D_{\Gamma_\infty,\Delta}$ all zeros remain exactly
\eqref{eq:surreal-zeros}.

\subsection{An uncountable hierarchy gives polynomial rigidity}
Let
\begin{equation}\label{eq:Gamma-omegaone}
 \Gamma_{\omega_1}=\bigoplus_{\alpha<\omega_1}\Q e_\alpha,
 \qquad e_\alpha=\omega^\alpha,
\end{equation}
again ordered by the largest nonzero index. Its cofinality is $\aleph_1$:
the basis is cofinal, whereas any countable family of finite supports has its
indices bounded below $\omega_1$. A larger basis element bounds the whole
family. Therefore
\begin{equation}\label{eq:uncountable-rigidity}
 \mathcal E_{\Gamma_{\omega_1}}
 =\mathcal T_{\Gamma_{\omega_1}}
 =K_{\Gamma_{\omega_1}}[Z].
\end{equation}
For the middle equality, a nonpolynomial restricted series would itself supply
a countable sequence of coefficient valuations tending cofinally to infinity.

\begin{center}\small
\begin{tabularx}{\textwidth}{@{}p{0.24\textwidth}cYYY@{}}
\toprule
Value group & Cofinality & Has an order unit & Nonpolynomial entire series & Nonconstant restricted units\\\midrule
$\Q$ or $\R$ & $\aleph_0$ & Every positive element & Yes & Yes\\[0.45em]
$\Q\epsilon\oplus\Q\eta$, $\epsilon\ll\eta$ & $\aleph_0$ & Yes; not every positive element & Yes & Yes, with the sharper gap test\\[0.45em]
$\Gamma_\infty$ & $\aleph_0$ & No & Yes & No\\[0.45em]
$\Gamma_{\omega_1}$ & $\aleph_1$ & No & No & No\\
\bottomrule
\end{tabularx}
\end{center}

\subsection{The whole surcomplex class}
\begin{corollary}[Whole-class ordinary-series rigidity]
\label{cor:whole-class}
Let $f=\sum_{n\ge0}a_nZ^n$ be an ordinary power series with surcomplex
coefficients. If its defining family is strongly Hahn-summable at every
$z\in\No[i]$, then $f$ is a polynomial.
\end{corollary}
\begin{proof}
The union of the normal-form supports of the countably many coefficients is
a set. The divisible additive span of this union together with $1$ is a
nonzero set-sized ordered subgroup $\Gamma\subset\No$, and all coefficients
lie in $K_\Gamma$.
The hypothesis makes $f$ entire over that field.

Every set of surreal numbers has a surreal upper bound: the set-sized cut
$\delta=\{\Gamma\cup\{0\}\mid\varnothing\}$ supplies a positive
$\delta\in\No$ larger than every member of $\Gamma$. Enlarge the group
to its span with $\delta$. If $f$ were nonpolynomial,
Theorem~\ref{thm:extension} would prohibit strong evaluation at $t^{-\delta}$,
contradicting the hypothesis.
\end{proof}

The repository already distinguishes such all-scale rigidity from the
existence of nonpolynomial series in fixed workspaces \cite{RepoMap,RepoRankOne}.
Here it is a consequence of the exact extension theorem. It does not imply
that every class function on $\No[i]$ is polynomial. It says nothing of that
kind about Gonshor's surreal exponential, locally represented function
classes, or series with a different indexing and summation convention.

\section{What the results do and do not establish}\label{sec:scope}

\subsection{Three independent distinctions}
The first distinction is between strong Hahn summation and intrinsic valuation
convergence. Positive well-ordered support suffices for the former; cofinal
progress through every group threshold is needed for the latter. On a single
disk they differ. For a series strongly evaluable at \emph{every} point,
Theorem~\ref{thm:criterion} proves that the distinction disappears within
its original field.

The second distinction is between countable cofinality and an order unit.
Countable cofinality supplies a sequence of ever larger scales, possibly
moving through infinitely many Archimedean classes. An order unit supplies
all those bounds by the integer multiples of one element. Entire functions
need only the first. A nonconstant restricted inverse requires the second.
The group $\Gamma_\infty$ separates them.

The third distinction is between an intrinsic topology in one workspace and
the fine topology tested against every surreal scale. The same countable
sequence may converge intrinsically in $K_\Gamma$ but fail to converge after
a noncofinal enlargement. The existence of an algebraic Hahn value can
survive even that failure. None of our proofs identifies these topologies.

\subsection{Why the positive-support preparation is indispensable}
One might try to derive everything from ordinary Tate-algebra preparation by
replacing real-valued norms with values in $\Gamma$. Example~\ref{ex:ranktwo-unit}
shows why that is unsafe. The standard iteration may improve the error by
$\epsilon$ repeatedly while remaining forever below $\eta$. Its purported
convergence is not proved.

Our preparation is instead an identity in the explicitly named ring
$\mathcal H_\Gamma$. The support certificate $S^*$ supplies every infinite
operation. On an individual disk its unit need not lie in
$\mathcal T_\Gamma$. In the global factorization proof, a quotient is shown
strongly evaluable at \emph{all} scales before Theorem~\ref{thm:criterion}
upgrades it to an entire valuation-convergent series. That ordering of the
arguments avoids a circular or false local inversion assumption.

\subsection{A useful topological interpretation}
The valuation topology of $K_\Gamma$ has a countable neighborhood base at zero
exactly when $\cf(\Gamma)=\aleph_0$. One direction uses a cofinal sequence
of valuation thresholds. Conversely, from a countable neighborhood base choose
one basic valuation ball inside each member. These selected thresholds must
be cofinal, since the base refines every other valuation ball.

A nonzero topologically nilpotent element exists exactly when $\Gamma$ has
an order unit, because $v(x^n)=nv(x)$. Hence the two analytic existence
criteria can be read as follows: nonpolynomial entire series are tied to a
countable valuation base, while nonconstant restricted units are tied to a
single topologically nilpotent scale. The latter distinction is consistent
with the higher-rank topological framework discussed by Conrad \cite{Conrad}.
We do not claim these underlying topological facts themselves as new.

\subsection{Limitations of the result}
Our exact extension domain keeps the coefficient field $\C$ fixed and uses
canonical ordered-group inclusions of full Hahn fields. It is not a theorem
about arbitrary extensions of incomplete subfields, arbitrary changes of
residue field, or arbitrary analytic sheaves. Algebraic closedness of the
Hahn workspaces is used essentially to conclude that all zeros of a prepared
polynomial already exist in the original workspace.

All entire functions studied here are represented by one ordinary formal
power series at the origin. Translation preserves this class, but that does
not identify it with every possible locally analytic or surreal-differentiable
function. In particular, piecewise functions on fine-topological domains and
special functions defined by other support conventions are outside the stated
hypotheses. The absence of a nonpolynomial whole-class series is therefore
not a general no-go theorem for surreal analysis.

The conclusions about finite zero sets, value distribution, and genus-zero
products parallel established non-Archimedean theory. Their inclusion is part of a complete
arbitrary-rank development, not a claim to have discovered canonical products
or Newton polygons.

\section{Proof audit, reproducibility, and further research}\label{sec:audit}

\subsection{Dependency audit}
The logical structure is deliberately separated into small parts.
The standard support lemmas give integral evaluation and support preparation.
Independently, the descending-subsequence argument gives the global criterion
and cofinality dichotomy. The restricted-unit theorem uses Gauss
multiplicativity and a convex coarsening; it does not use zero counting or
canonical products. The extension-domain theorem uses the global definition
and integral evaluation; it does not use factorization. Only zero
conservation imports algebraic closedness and preparation. The divisor
classification is downstream of these local tools and the elementary entire
algebra.

Thus the proposed principal classifications are not consequences of an
unproved global factorization assertion. Conversely, the factorization proof
never assumes that an arbitrary zero-free restricted series is restricted-invertible;
that is exactly the assertion that can fail in higher rank.

\subsection{The exact-arithmetic companion}
The archive includes \code{code/verify.py}, a standalone Python 3 program
using only the standard library. Exponents are tuples of exact rational
numbers, ordered by the largest nonzero coordinate. Coefficients and
polynomial division are exact rational arithmetic. No floating-point fitting
or numerical root search is used.

The recorded run reports \textbf{664 passing finite checks}. It verifies
finite truncations of \eqref{eq:infinite-product}, their exact zeros and
simple-root derivatives, elementary-symmetric leading exponents, finite
Newton closed- and open-ball counts, the descending exponent obstruction
after adjoining a dominating scale, the hidden-tail examples, and the
preparation recursion through order eight for \eqref{eq:cubic}.
The detailed machine-readable record is \code{data/verification.json};
\code{data/verification.txt} is its short summary.

These checks can detect sign mistakes, reversed order conventions, incorrect
root-count endpoints, and wrong finite preparation coefficients. They cannot
prove well-ordering of arbitrary supports, the cofinality dichotomy, necessity
in the unit criterion, or infinite product identities. Those assertions rely
on the mathematical proofs above. No Lean verification is claimed.

\subsection{Literature and novelty audit}
The primary-source review used Poonen's Hahn-field construction and algebraic
closedness result, Cherry's rank-one lectures and factorization paper,
Conrad's higher-rank valuation examples, and the surreal normal-form
literature \cite{Poonen,CherryLectures,CherryFactorization,Conrad,BournezGuilmant}.
The exact cofinality classification for global \emph{strong} evaluation and
the universal scalar-extension domain were not found in those sources. The
restricted-unit theorem is closely related to established topological
nilpotence distinctions; its precise necessity-and-sufficiency statement for
\eqref{eq:T}, including groups with no order unit, is the formulation proposed
here.

The search is not exhaustive. A proof in the present article is a mathematical
claim supported by its argument; a first-publication claim would require a
more extensive bibliographic and expert review. We therefore describe the
principal results as proposed new arbitrary-rank formulations and deductions,
not as certified historical firsts. The dated repository audit and the external
references are included to make that distinction inspectable.

\subsection{Further questions suggested by the proofs}
The single-variable exact extension domain is unusually uniform. In several
variables a direct assertion that every nonpolynomial entire series has the
same product domain is false: a function independent of one coordinate places
no summability restriction on that coordinate. The finite-variable convergence
and unit criteria in Appendix~\ref{app:multivariate} suggest that the correct
replacement should depend on the unbounded degree directions of the support.
A useful next problem is to characterize the maximal strong extension domain
by that asymptotic degree geometry, including its boundary where higher Hahn
tails may matter.

A second question is how much of the divisor theory survives for incomplete
surreal subfields or for support-restricted Hahn subfields. Our proof isolates
two requirements: closure under the positive-support preparation construction,
and enough algebraic roots to split its finite prepared polynomials. A
classification of natural surreal subfields satisfying those requirements
would turn the present full-Hahn results into a more computationally selective
theory.

For formalization, the first useful target is the global criterion: its
nontrivial implication is a short ordered-group argument plus the definition
of strong summability. The next target is the coarsened polynomial argument
for units. The support preparation should be formalized only after the
finite-word Neumann lemma and arbitrary regrouping are available. This is a
dependency plan, not a claim that the proofs have already been checked in Lean.

\section{Conclusion}
The threshold for nonpolynomial entire behavior in a Hahn--surcomplex
workspace is countable cofinality, not rank one. The threshold for a nonconstant
restricted inverse is stronger: a positive coefficient gap must generate a
cofinal sequence of multiples. These conditions can separate in an explicit
infinite-rank subgroup of the surreal numbers.

After enlarging the value group, every nonpolynomial entire series has the
same exact strong domain: the valuation ring of the convex hull of its
original scales. Cofinality is precisely what preserves entireness, while
support preparation preserves every existing zero and prevents new ones
throughout the surviving domain. The canonical product theorem then shows
that coefficients, zeros and the extension boundary encode the same cofinal
scale structure. This gives a unified explanation of both nonpolynomial
fixed-workspace functions and polynomial rigidity on the whole surcomplex
class, without confusing Hahn identities with convergence in the fine
surreal topology.

\appendix
\section{Finite-variable convergence and unit criteria}\label{app:multivariate}
Let $d\ge1$ be finite, write
$\alpha=(\alpha_1,\ldots,\alpha_d)\in\N^d$ and
$|\alpha|=\alpha_1+\cdots+\alpha_d$. Define
\[
 \mathcal T_{\Gamma,d}
 =\left\{\sum_\alpha a_\alpha X^\alpha:
              v(a_\alpha)\to+\infty\text{ as }|\alpha|\to\infty\right\}.
\]
Define $\mathcal E_{\Gamma,d}$ by strong summability of the \emph{whole}
family $(a_\alpha x^\alpha)_{\alpha\in\N^d}$ for every
$x\in K_\Gamma^d$. Summing first along one coordinate and then another is
not the definition.

\begin{theorem}\label{thm:multi-criterion}
For a finite-variable formal series $F=\sum_\alpha a_\alpha X^\alpha$,
\[
 F\in\mathcal E_{\Gamma,d}
 \quad\Longleftrightarrow\quad
 \frac{v(a_\alpha)}{|\alpha|}\to+\infty
 \quad(|\alpha|\to\infty,\ |\alpha|>0).
\]
It is enough to test strong summability at the diagonal monomial points
$(t^s,\ldots,t^s)$ for all $s\in\Gamma$.
Nonpolynomial members exist exactly when $\cf(\Gamma)=\aleph_0$.
\end{theorem}
\begin{proof}
At a diagonal monomial point the summand values are
$v(a_\alpha)+|\alpha|s$. There are only finitely many multi-indices of any
bounded total degree. If these values do not tend cofinally to infinity,
strong summability gives a common lower bound and an infinite bounded-above
subfamily with strictly increasing total degrees. Subtracting
$|\alpha|\delta$, for a positive $\delta$ larger than the bounded interval's
length, gives the same strictly descending leading-exponent sequence as in
Theorem~\ref{thm:criterion}. Thus diagonal strong summability implies the
slope criterion.

Conversely, for a point with all coordinates nonzero, set
$s=\min_i v(x_i)$. Then
\[
 v(a_\alpha x^\alpha)
 \ge v(a_\alpha)+|\alpha|s\to+\infty.
\]
Because bounded total degrees contain finitely many indices,
Lemma~\ref{lem:cofinal-strong} applies after any enumeration respecting
total degree. If some coordinates are zero, discard the vanishing terms
and apply the same argument to the remaining coordinates. The cofinality
assertion follows exactly as in one variable; for existence one may use an
entire nonpolynomial in $X_1$ alone.
\end{proof}

\begin{theorem}\label{thm:multi-unit}
A nonconstant $F\in\mathcal T_{\Gamma,d}$ is a unit in that algebra if and
only if $a_0\ne0$ and
\[
 \min_{\alpha\ne0,\ a_\alpha\ne0}v(a_\alpha/a_0)
\]
is a positive order unit of $\Gamma$.
\end{theorem}
\begin{proof}
Regrouping into $\C[X_1,\ldots,X_d]((t^\Gamma))$ gives multiplicativity of
the Gauss valuation. If $FG=1$, normalize both constant coefficients to $1$.
Their Gauss values are nonpositive and sum to zero, so their initial
polynomials multiply to $1$ and are constant. All nonconstant coefficient
values are positive. If their minimum $\delta$ is not an order unit,
coarsen by the proper convex subgroup it generates. Restrictedness reduces
both series to finite polynomials over the coarsened residue field, and
at least one nonconstant coefficient of $F$ survives. This contradicts
invertibility of its reduction. Conversely an order-unit gap makes the
geometric inverse restricted, with precisely the tail estimates in
Theorem~\ref{thm:units}. Finite-variable convolution has finitely many terms
at each multi-index, so those estimates are unchanged.
\end{proof}

The finite-dimensionality hypothesis has been used in the finiteness of
bounded-degree index sets. These statements do not assert the corresponding
results for infinitely many variables. Nor do they extend the
one-variable divisor classification to higher-dimensional zero sets.

\clearpage
\phantomsection
\addcontentsline{toc}{section}{References}
\begingroup
\small\raggedright
\begin{thebibliography}{99}

\bibitem{Neumann}
B.~H. Neumann,
\emph{On ordered division rings},
Transactions of the American Mathematical Society \textbf{66} (1949),
202--252.

\bibitem{Poonen}
B.~Poonen,
\emph{Maximally complete fields},
L'Enseignement Math\'ematique (2) \textbf{39} (1993), 87--106.
In particular Section~3 and Corollary~4.
\url{https://doi.org/10.5169/seals-60414}.
Accessible scan:
\url{https://virtualmath1.stanford.edu/~conrad/Perfseminar/refs/poonencomplete.pdf}.

\bibitem{Gonshor}
H.~Gonshor,
\emph{An Introduction to the Theory of Surreal Numbers},
London Mathematical Society Lecture Note Series 110,
Cambridge University Press, 1986.

\bibitem{BournezGuilmant}
O.~Bournez and Q.~Guilmant,
\emph{Surreal fields stable under exponential and logarithmic functions},
2022, arXiv:2201.08199.
\url{https://arxiv.org/abs/2201.08199}.

\bibitem{CherryLectures}
W.~Cherry,
\emph{Lectures on Non-Archimedean Function Theory},
2009, arXiv:0909.4509.
In particular the preparation and product results in Lectures~2--3.
\url{https://arxiv.org/abs/0909.4509}.

\bibitem{CherryFactorization}
W.~Cherry,
\emph{Existence of GCD's and Factorization in Rings of Non-Archimedean
Entire Functions},
Contemporary Mathematics \textbf{551} (2011), 57--69;
arXiv:1007.0984v2.
\url{https://arxiv.org/abs/1007.0984}.

\bibitem{Conrad}
B.~Conrad,
\emph{Lecture 6: More constructions with Huber rings},
Stanford Number Theory Learning Seminar on perfectoid spaces,
31 October 2014, informal lecture notes.
In particular Section~6.2 and Example~6.2.1.
\url{https://virtualmath1.stanford.edu/~conrad/Perfseminar/Notes/L6.pdf}.
Seminar attribution:
\url{https://virtualmath1.stanford.edu/~conrad/Perfseminar/}.

\bibitem{RepoMap}
V.~Reshetnikov, maintainer,
\emph{Surreal repository: the surreal and surcomplex reports},
\code{docs/README.md}; AI-assisted research collection,
snapshot \code{aa846271b4dc}, inspected 21 September 2026.
\url{https://github.com/VladimirReshetnikov/Surreal/blob/aa846271b4dcae2c055b216126a87210292ec19b/docs/README.md}.

\bibitem{RepoRankOne}
\emph{From Surcomplex Hahn Fields to Berkovich Geometry:
Rank-one convergence, Tate algebras, analytic annuli, and an entire-function
case study}, AI-assisted research report in the Surreal repository,
\code{docs/surcomplex/rank-one-berkovich/article.tex},
snapshot \code{aa846271b4dc}, 21 September 2026.
The existing report already contains rank-one zero-free rigidity and a
canonical product theorem.
\url{https://github.com/VladimirReshetnikov/Surreal/blob/aa846271b4dcae2c055b216126a87210292ec19b/docs/surcomplex/rank-one-berkovich/article.tex}.

\end{thebibliography}
\endgroup
\end{document}
